\section{Mixture-of-Experts: Theory and Implementation}
\label{app:moe}

This appendix documents the Mixture-of-Experts (MoE) contract implemented by
\texttt{src/model/moe/} and declared in
\texttt{src/schema/\_model/\_moe.yaml}. The schema is the executable source of
truth: every key discussed here exists in the schema, is materialised in the
nested \texttt{FrankensteinModelConfig} \texttt{moe} block, and is exercised by
the test suite. The first two sections are the on-ramp: they explain what an
MoE layer does---one equation, one analogy, one worked example---and then the
two failure modes (routing collapse and the hot-expert straggler) that force
most of the design. The rest is the reference: router scoring, routing
strategy, load balancing, experts, layer layout, capacity and monitoring,
closing with the schema-contract and preset tables. A reader who finishes
Section~\ref{app:moe-failures} should be able to open any later section and
know which of the two failures, or which of the seven design questions, it is
answering.

\subsection{Overview: what a Mixture-of-Experts layer is}
\label{app:moe-overview}

The feed-forward half of a transformer block is \emph{dense}: one network,
applied identically to every token, no matter what the token is. That is a
clinic with a single physician who must treat every complaint. An MoE layer
turns the clinic into a roster of $E$ specialists with a receptionist: the
\emph{router} reads each token and sends it to the $k$ most relevant
specialists, the \emph{experts} $\{E_e\}_{e=1}^{E}$. Formally,
\begin{equation}
  y_t = g \sum_{j=1}^{k} \hat w_{t,j}\, E_{i_{t,j}}(x_t) + y^{\text{shared}}_t ,
  \label{eq:moe-layer}
\end{equation}
read from the inside out: expert $i_{t,j}$ processes the token and produces a
vector, the mixing weight $\hat w_{t,j}$ scales it, the $k$ contributions are
summed, the routed scaling factor $g$ restores the branch's magnitude
(Section~\ref{app:moe-router}), and $y^{\text{shared}}_t$ is the optional
always-on shared-expert branch (Section~\ref{app:moe-experts}). Only the
selected experts run, which is the whole bargain: total parameters grow with
$E$ while per-token compute stays proportional to $k$
\citep{shazeer2017moe,fedus2021switch}.

The bargain in numbers: with $E = 8$ experts of the dense width and $k = 2$,
the layer stores $8\times$ the feed-forward parameters but each token is
processed by only $2$ experts---$25\%$ of the pool---plus a router of $d
\times E$ weights, which is negligible next to any expert. Memory grows with
$E$; FLOPs grow with $k$. Those are two independent budgets, and most of the
schema is the machinery for spending them well.

Seven questions span the design space, and Table~\ref{tab:moe-questions} maps
each to its knobs and its section: (i) how logits become weights, (ii) who
selects the experts, (iii) how load is kept balanced, (iv) what the experts
look like, (v) which layers are MoE at all, (vi) what happens when an expert
overflows, and (vii) how all of it is observed during training.

\begin{table}[H]
\centering
\caption{The seven questions of an MoE layer, the knobs that answer them and
the section that develops each answer. Every \texttt{model.moe.*} key of the
schema contract (Table~\ref{tab:moe-schema}) belongs to exactly one row.}
\label{tab:moe-questions}
\small
\begin{tabular}{@{}L{2.2cm}L{4.8cm}L{4.6cm}L{1.6cm}@{}}
\toprule
\textbf{Axis} & \textbf{Question} & \textbf{Knobs} & \textbf{Section} \\
\midrule
Scoring & how much weight does each selected expert get? & \texttt{router.*} & \ref{app:moe-router} \\
Routing & who picks whom? & \texttt{routing.type}, \texttt{routing.hash\_*} & \ref{app:moe-routing} \\
Balancing & how do we keep every expert employed? & \texttt{balancing.*} & \ref{app:moe-balancing} \\
Experts & how many, how wide, shared or latent? & \texttt{num\_experts}, \texttt{expert\_hidden\_size}, \texttt{shared.*}, \texttt{latent.*}, \texttt{experts.*} & \ref{app:moe-experts} \\
Layout & which layers are MoE at all? & \texttt{layout.*} & \ref{app:moe-layout} \\
Capacity & what happens when an expert overflows? & \texttt{routing.capacity\_factor} & \ref{app:moe-capacity} \\
Monitoring & how is all of it observed? & logged scalars & \ref{app:moe-monitoring} \\
\bottomrule
\end{tabular}
\end{table}

\subsection{The two failure modes: collapse and the straggler}
\label{app:moe-failures}

Four of the seven questions exist because of two failure modes. Both follow
from one design decision---the router is learned and free to concentrate---and
every mechanism in the rest of this appendix (auxiliary loss, expert bias,
expert choice, hash tables, capacity factor) is an answer to one of them.

\paragraph{Collapse.}
Routing decisions are uncoordinated: each token independently picks its own
$k$ favourites, and nothing stops every token from picking the same expert.
The loop is self-reinforcing. An expert that starts slightly ahead receives
more tokens; more tokens mean more gradient, so it improves faster than the
others; improving faster, it wins even more tokens on the next batch. Run the
loop forward and a layer built with $E$ experts routes through two or three
of them, and the rest never train at all. Switch Transformers observed whole
layers collapsing to a single expert \citep{fedus2021switch}. This is the
default trajectory of an unbalanced MoE, not a rare accident, and it is why
Section~\ref{app:moe-balancing} exists: the auxiliary loss and the expert bias
are two independent ways to break the loop, expert choice
(Section~\ref{app:moe-routing}) removes it by construction, and hash layers
remove it by fixing the selection outright.

\paragraph{The straggler.}
Even a layer that balances well on average fluctuates per step: one batch may
send a third of its tokens to one expert. Two costs follow. The step cannot
finish until the busiest expert finishes, so a single hot expert sets the time
of everything, and in distributed training---where different experts live on
different devices---the imbalance becomes rows of idle devices waiting for
one. And an expert that receives a burst needs buffer memory proportional to
the burst. The capacity factor (Section~\ref{app:moe-capacity}) bounds both
by capping how many assignments an expert must honour. The gauge for both
failure modes is MaxVio (Section~\ref{app:moe-balancing}): collapse and
stragglers both read as the busiest expert pulling away from the mean.

\subsection{Router scoring: from logits to weights}
\label{app:moe-router}

The router is one linear layer, $z_t = W_r x_t \in \mathbb{R}^{E}$ with
$W_r \in \mathbb{R}^{E \times d}$: one logit per expert. A gating function
converts the $E$ logits into $E$ scores, and the top-$k$ scores select the
experts. It helps to separate the two jobs of the gating function: it
\emph{ranks} (which experts are selected) and it \emph{weights} (how much each
selected expert contributes). The implementation supports the three gating
functions used by current open models:
\begin{align}
  \text{softmax}:&\quad s_{t,e} = \frac{e^{z_{t,e}}}{\sum_{e'} e^{z_{t,e'}}}, \\
  \text{sigmoid}:&\quad s_{t,e} = \sigma(z_{t,e}), \\
  \text{sqrt-softplus}:&\quad s_{t,e} = \sqrt{\operatorname{softplus}(z_{t,e})}.
\end{align}
All three are monotone in $z_{t,e}$---a larger logit always yields a larger
score---and monotone functions never reorder. So for a given $z_t$ all three
select \emph{exactly the same experts}; the choice of gating function is a
choice about the weights, not about the selection. In words:
\begin{itemize}[leftmargin=1.4em]
  \item \textbf{softmax} makes experts compete for a fixed pie. The scores sum
        to one, so expert $e$'s score depends on every other logit through the
        denominator: raising one logit lowers all the others. As logits grow,
        the distribution sharpens towards winner-takes-all.
  \item \textbf{sigmoid} judges each expert in isolation. $\sigma(z_{t,e})$
        depends only on its own logit, so there is no competition: all experts
        can score low together (the selected weights sum to less than one and
        the routed branch shrinks) or high together. The renormalisation and
        scaling factor below exist largely to manage exactly this freedom
        \citep{deepseek2024v3}.
  \item \textbf{sqrt-softplus} sits between the two: like sigmoid it is
        per-expert, but $\sqrt{\operatorname{softplus}(z)}$ decays more gently
        on the negative side and never saturates at exactly zero
        \citep{deepseek2026v4}.
\end{itemize}
\texttt{router.scoring} selects the function (default \texttt{softmax}).

\paragraph{Output scale.}
Two knobs control the magnitude of the routed branch, and they only make
sense as a pair. With \texttt{router.renormalize} (default true), the
selected scores are rescaled to sum to one,
\begin{equation}
  \hat w_{t,j} = \frac{s_{t,i_{t,j}}}{\sum_{l=1}^{k} s_{t,i_{t,l}}},
\end{equation}
so the routed branch becomes a weighted \emph{average} of $k$ expert outputs
rather than a sum. Averaging shrinks. Take $k = 4$ experts whose outputs each
have root-mean-square $1$ and are uncorrelated, and near-equal weights of
$\approx 1/4$: the variance of the average is
$\tfrac{1}{4^2}(1+1+1+1) = \tfrac{1}{4}$, so its RMS is $1/\sqrt{k} = 0.5$.
The average is half the size of any single expert's answer, and the shrinkage
worsens as $k$ grows. Renormalisation still pays for itself---it removes the
token-to-token variation of the top-$k$ probability mass and lowered
validation loss in the granularity study of \citet{krajewski2025scaling}---so
the routed scaling factor $g$ (\texttt{router.routed\_scaling\_factor},
default $1$) is there to undo the shrinkage: $g \approx \sqrt{k}$ restores
the RMS of a single expert. The published values match the derivation. Kimi
K2 uses $2.827 \approx \sqrt{8}$; DeepSeek-V3 uses $g = 2.5$ at $k = 8$,
slightly below $\sqrt{8} \approx 2.83$ because real expert outputs are partly
correlated and cancel less than independent ones
\citep{deepseek2024v3}; and Kimi K3 uses $g = 1.0$ at $k = 16$ with two
shared experts carrying part of the branch \citep{kimiteam2026k3}.

\paragraph{Pre-softmax.}
With \texttt{router.pre\_softmax} the softmax runs over all $E$ experts
\emph{before} the top-$k$ cut, so the kept weights are raw probabilities from
the full distribution and are not renormalised. To see why softmax routing at
$k{=}1$ requires this, trace the default path step by step: select first,
then softmax over the \emph{selected} logits only. With one selected logit
that softmax is $e^{z}/e^{z} = 1$---always, regardless of how well the expert
matches the token. A weight that is constantly one has derivative zero, so
the router would receive no gradient and could never learn which expert to
prefer: selection would be frozen at its initialisation. Pre-softmax repairs
it: the kept weight is the top probability of the full softmax, a genuine
number in $(0,1)$ that moves whenever any logit moves, so the gradient flows.
The configuration layer detects that combination, enables
\texttt{pre\_softmax} automatically, and emits a warning instead of rejecting
legacy configs.

\paragraph{Sharpness and precision.}
Three mechanisms control how peaked, and how numerically trustworthy, the
router is. \emph{Input jitter} (\texttt{router.jitter\_eps}) multiplies the
token by $x \cdot (1 + \epsilon(2u-1))$ with $u$ uniform before the
projection: random multiplicative noise that penalises routing decisions
built on fragile directions of $x_t$, a regulariser Switch was trained with
\citep{fedus2021switch}. \emph{Router precision} (\texttt{router.dtype})
defaults to \texttt{float32}; the router is tiny ($d \times E$ parameters),
so full precision costs nothing, and what it buys is stability---a bf16
router diverged in Switch Transformers, and in low precision tied top-$k$
logits are broken arbitrarily \citep{fedus2021switch}. Finally, the
\emph{z-loss} of ST-MoE penalises logit drift. The log-sum-exp is a smooth
maximum of the logits (large whenever \emph{any} logit is large), so
\begin{equation}
  \mathcal{L}_z = \beta\,\frac{1}{T}\sum_{t=1}^{T}
  \Bigl(\log \sum_{e=1}^{E} e^{z_{t,e}}\Bigr)^{2},
  \label{eq:zloss}
\end{equation}
grows quadratically whenever any router logit grows in magnitude. Minimising
it keeps the logits small, the softmax unsaturated and the router from
becoming overconfident \citep{zoph2022stmoe}. Here $\beta$ is
\texttt{router.z\_loss\_coeff} (default $0$; ST-MoE uses $10^{-3}$); the loss
is returned as \texttt{moe\_router\_z\_loss} and added to the task loss.

\subsection{Routing strategy: who selects, and how}
\label{app:moe-routing}

Section~\ref{app:moe-router} fixed \emph{how much} each chosen expert weighs;
this one fixes \emph{who chooses whom}. The literature has three answers, and
each has a structural consequence.

\paragraph{Token choice.}
The default (\texttt{routing.type: token\_choice}): each token keeps its own
$k$ best experts under $s_t + b$, where $b$ is the balancing bias of
Section~\ref{app:moe-balancing}. The consequence is that nobody guarantees any
expert work: an expert can win every token or none, which is exactly the
collapse loop of Section~\ref{app:moe-failures}. Token choice therefore ships
with a balancing mechanism. It is the strategy trained by Megatron and used
by most open models \citep{fedus2021switch,deepseek2024v3}.

\paragraph{Expert choice.}
Flip the direction of the decision: with
\texttt{routing.type: expert\_choice}, each expert keeps its best $Tk/E$
tokens \citep{zhou2022expertchoice}. Every expert receives exactly the same
number of tokens, so loads are equal by construction and the balancing
machinery is not merely unnecessary but contradictory---the schema forbids
combining expert choice with an auxiliary loss or an expert bias. The price
is causality: whether a token is served now depends on the \emph{other}
tokens in the batch, including later ones. That is harmless in an encoder,
where the whole sequence is available at once, and fatal in autoregressive
decoding, where token $t$ must be routed before tokens $t{+}1, t{+}2,
\dots$ exist; training-time and inference-time routing would disagree by
construction. Hence the recommendation: encoder and training experiments
yes, autoregressive serving no.

\paragraph{Hash layers.}
Dispense with the learned decision altogether: the first
\texttt{routing.hash\_layers} MoE layers select experts from a fixed table
$\mathrm{tid} \rightarrow (e_1,\dots,e_k)$ \citep{roller2021hash}, as
DeepSeek-V4 does in its first three MoE layers \citep{deepseek2026v4}. The
observation behind it: in the first layers, tokens are still close to their
embeddings and the routing decision is mostly surface identity---which token
it is, not what it means---and a lookup table captures that. The table never
changes, so it cannot collapse. The learned router still produces the mixing
weights, so it keeps training through them; it just receives no gradient from
the selection itself. The table is built from a seeded random assignment
(\texttt{hash\_table: random}) or a round-robin cycle (\texttt{round\_robin});
Roller et al.\ found random and frequency-greedy tables reach similar
perplexity, with the greedy table better balanced \citep{roller2021hash}.
Because a fixed table can never be rebalanced, the schema forbids combining
it with \texttt{balancing.bias\_update}.

\subsection{Load balancing}
\label{app:moe-balancing}

Section~\ref{app:moe-failures} described the collapse loop: an expert
slightly ahead receives more tokens, trains faster, and pulls further ahead.
Two independent remedies are implemented; either works alone, and DeepSeek-V3
combines both \citep{deepseek2024v3}. They differ in \emph{how} they reach
the routing: the auxiliary loss works through the gradient, the bias works
around it.

\paragraph{Auxiliary loss.}
The obstacle to fixing imbalance directly is that the thing to fix---how many
times each expert was chosen---is a count, and counts have no derivative:
being chosen $500$ times instead of $499$ is not an operation backpropagation
can influence. The Switch surrogate \citep{fedus2021switch} multiplies each
expert's share of the assignments by its mean router probability,
\begin{equation}
  \mathcal{L}_{\text{aux}} = \alpha\,E \sum_{e=1}^{E} f_e P_e,
  \qquad
  f_e = \frac{1}{Tk}\sum_{t,j}\mathbb{1}[i_{t,j}=e],
  \qquad
  P_e = \frac{1}{T}\sum_{t} s_{t,e},
  \label{eq:aux}
\end{equation}
where $f_e$ (the fraction of the $Tk$ assignments) carries the imbalance but
no gradient, and $P_e$ (the mean score) carries gradient. The product is the
trick: lowering the loss means lowering $P_e$ precisely where $f_e$ is high,
so the optimiser leans on the router's probabilities for exactly the
overloaded experts, and future selections spread. Under softmax scoring the
scores sum to one per token and the minimum sits at the uniform
distribution, $f_e = k/E$ and $P_e = 1/E$; the prefactor $E$ then makes the
loss at the minimum independent of the expert count (exactly $\alpha$ at
$k{=}1$), so one coefficient means the same thing at $E{=}8$ and $E{=}256$.
The coefficient itself is a trade-off: too large and it fights the
language-model loss, too small and MaxVio stays high
\citep{wang2024auxfree}; the default is $0.01$, close to Switch and
DeepSeekMoE \citep{fedus2021switch,dai2024deepseekmoe}. The scope is
selected by \texttt{balancing.aux\_loss}: \texttt{batch} (one micro-batch,
the default), \texttt{sequence} (one loss per sequence, DeepSeek-V3's guard
against a single sequence piling onto a few experts---balance the batch and
one long document can still monopolise its favourite experts), or
\texttt{none}. Balancing per sequence or per micro-batch can prevent domain
experts from forming, whereas global-batch balancing leaves each device's
load uneven \citep{qiu2025demons}; the single-device runtime implements the
first two scopes and documents the global scope as future work.

\paragraph{Auxiliary-loss-free bias.}
The second remedy ignores gradients entirely and works like a thermostat:
measure the temperature, nudge the knob, repeat. With
\texttt{balancing.bias\_update} set to \texttt{sign} or
\texttt{centered\_sign}, each expert carries a bias $b_e$ added to the scores
\emph{before selection only---never to the mixing weights}---and moved once
per optimizer step:
\begin{equation}
  b_e \leftarrow b_e + \gamma\,\operatorname{sign}\!\left(\bar c - c_e\right),
  \qquad
  \bar c = \frac{1}{E}\sum_{e'} c_{e'},
  \label{eq:bias}
\end{equation}
where $c_e$ counts the step's assignments and $\gamma$ is
\texttt{balancing.bias\_update\_rate} (default $10^{-3}$). An expert above
the mean load ($c_e > \bar c$) has its bias lowered by $\gamma$, so it wins
fewer tokens on the next step; an expert below the mean has its bias raised.
Because the bias never touches the mixing weights, it redirects traffic
without changing what the selected experts contribute---which is why it
costs no language-model loss, unlike the auxiliary loss
\citep{wang2024auxfree}.

Why the schema restricts the bias to sigmoid or sqrt-softplus scoring: under
softmax every score depends on every logit through the denominator, so
shifting one expert's logit rescales all $E$ scores at once and the bias
would steer the selection \emph{and} distort every mixing weight at the same
time. Under sigmoid or sqrt-softplus each score is a function of its own
logit alone, so $b_e$ moves expert $e$ in the ranking and touches nothing
else---steering stays pure.

The \texttt{centered\_sign} variant subtracts the mean update
($\tfrac{1}{E}\sum_{e}\operatorname{sign}(\bar c - c_e) = 0$ in exact
arithmetic, not necessarily in reduced precision) so the bias mean does not
drift, which matters when the bias is stored in reduced precision
\citep{wang2024auxfree}. The bias is a non-persistent fp32 buffer: it is
training state, never an inference parameter, and \texttt{eval()} freezes it.
The update is applied by \texttt{update\_moe\_expert\_bias()} after
\texttt{optimizer.step()} in the trainer.

\paragraph{MaxVio.}
Balance is read from the maximal violation,
\begin{equation}
  \mathrm{MaxVio} = \frac{\max_e c_e - \bar c}{\bar c},
\end{equation}
a single number per layer, logged as \texttt{moe/max\_vio}. How to read it:
$0$ is perfectly even; $0.2$--$0.5$ is ordinary fluctuation; $1$ means the
busiest expert receives twice the mean load \citep{wang2024auxfree}. The
absolute value matters less than the trend---a MaxVio that climbs over
steps is the collapse loop of Section~\ref{app:moe-failures} running in
real time.

\subsection{Experts: count, width, latent space, sharing}
\label{app:moe-experts}

\paragraph{Count and width.}
The two budgets of Section~\ref{app:moe-overview} are spent separately here:
\texttt{num\_experts} $\times$ \texttt{expert\_hidden\_size} is (roughly) the
total routed parameters---what memory pays---while \texttt{top\_k} $\times$
\texttt{expert\_hidden\_size} is the per-token compute---what the user pays.
\texttt{expert\_hidden\_size} sets each expert's intermediate width and falls
back to \texttt{ffn.hidden\_size} when unset. Keeping the two widths separate
is essential for real architectures: DeepSeek-V3 has a dense width of
$18432$ but experts of width $2048$ \citep{deepseek2024v3}---its dense layers
are nine times wider than its experts. Granularity is the refinement of the
same trade: at a fixed parameter budget, ``$16$ experts of width $4096$''
and ``$64$ experts of width $1024$'' hold the same total. The finer split
lets each expert specialise more narrowly, which improved loss up to $128$
experts of width $d/4$ with top-$8$ in one study \citep{liew2026principled};
but narrower experts also see fewer tokens each per step, and once training
throughput is priced in, the same total and active parameters are predicted
to favour a \emph{less} sparse configuration \citep{krajewski2025scaling}.
The honest summary: finer helps at fixed FLOPs, and the evidence at fixed
wall-clock is mixed.

\paragraph{Latent MoE.}
Ship a compressed copy instead of the whole token. LatentMoE wraps the
experts between a down- and an up-projection: each token is projected from
$d$ to a latent width $\ell$ before dispatch, and the combined result is
projected back after \citep{elango2026latentmoe},
\begin{equation}
  \ell_t = W_\downarrow x_t,\qquad
  y_t = W_\uparrow \sum_{j=1}^{k} \hat w_{t,j} E_{i_{t,j}}(\ell_t).
  \label{eq:latent}
\end{equation}
Each of the $k$ dispatched copies now carries $\ell$ numbers instead of $d$.
With compression ratio $\alpha = d/\ell$, moving and processing one copy
costs $\alpha$ times less, so the same expert parameters, compute and
traffic afford $\alpha$ times more experts and a larger top-$k$. The bet is
that experts can work on a lossy sketch of the token; NVIDIA reports it
holds up to $\alpha \leqslant 4$ ($+5.65$ MMLU-Pro points at 95B total with
$\alpha = 4$) \citep{elango2026latentmoe}. Kimi K3 adds an RMSNorm before
the up-projection to stop routed-branch activation growth at 2.8T
parameters, exposed as \texttt{latent.norm\_before\_up\_proj}
\citep{kimiteam2026k3}.

\paragraph{Shared experts.}
A routed layer has no guarantee that common knowledge gets represented: each
expert learns whatever its visitors need, so $E$ experts may each re-learn
the same syntax and punctuation. \texttt{shared.count} always-on experts run
outside the router on \emph{every} token---the general practitioners of the
clinic in Section~\ref{app:moe-overview}---absorbing the knowledge every
token needs and freeing the routed specialists to spend their capacity on
what differs across tokens \citep{dai2024deepseekmoe}. Their width is
\texttt{shared.hidden\_size} (default: the routed width). They are added at
weight one, or multiplied by a learned per-token gate
$\sigma(w_s^{\top} x_t) \in (0,1)$ when \texttt{shared.gate: sigmoid}
(Qwen3-Next) \citep{qwen2025qwen3}; the fixed-weight form matches
DeepSeek-V3 \citep{deepseek2024v3}. They are not free: a shared expert runs
on every token, so per-token compute grows by its width, and measured
quality gains are small and inconsistent \citep{li2026slicing}. Add one when
you would otherwise raise \texttt{top\_k}, not on top of it.

\paragraph{Activation clamp.}
A rare but real pathology: individual experts whose output norms grow while
the loss stays flat---Step 3.5 Flash found such outliers, and their growth
was contained by clipping at unchanged loss \citep{huang2026step35}.
\texttt{experts.activation\_clamp} bounds the gated intermediate,
\begin{equation}
  \mathrm{GatedFFN}(x) = W_{\text{down}}\!\left(
  \phi\bigl(\operatorname{clamp}(xW_{\text{gate}}, -c, c)\bigr) \odot (xW_{\text{up}})
  \right),
\end{equation}
with $c$ = \texttt{activation\_clamp} ($0$ disables it). DeepSeek-V4.1 uses
$c = 10$ and Step 5 Preview $c = 7$ \citep{deepseek2026v4flash}. Treat it as
insurance against rogue experts, not as a default: it exists so that the
monitoring signal of Section~\ref{app:moe-monitoring} (expert-norm ratio
rising) has a cheap remedy. The exact SiTU-GLU smooth cap used by Kimi K3 is
approximated here by a gate-input clamp and documented as a deviation
\citep{kimiteam2026k3}.

\subsection{Layer layout}
\label{app:moe-layout}

Not every layer needs to be an MoE layer. The first layers of a transformer
compute generic features that every token needs regardless of content;
specialisation pays later. \texttt{layout.dense\_layers} keeps the first $N$
layers dense (DeepSeek-V3 and GLM keep three), and
\texttt{layout.every\_n\_layers} interleaves dense and MoE layers (Llama 4
Maverick makes every second layer an MoE layer)
\citep{deepseek2024v3,meta2025llama4}. Dense layers use
\texttt{ffn.hidden\_size}; MoE layers use \texttt{expert\_hidden\_size}. The
model config exposes \texttt{is\_moe\_layer(i)}, \texttt{moe\_layer\_mask()}
and \texttt{moe\_layer\_count()}; the encoder, decoder, ViT and audio stacks
build a routed or dense FFN per physical layer accordingly. Leading dense
layers reduce the number of MoE layers at little cost at matched FLOPs
\citep{krajewski2025scaling}, but they do not by themselves fix early-layer
imbalance: the first MoE layers after the prefix are still the most
collapse-prone and remain the ones to monitor.

\subsection{Capacity and token dropping}
\label{app:moe-capacity}

By default routing is dropless: an expert accepts every token the router
sends it (\texttt{routing.capacity\_factor: 0}), and per-step loads fluctuate
freely. A positive factor $\kappa$ caps each expert at
\begin{equation}
  C = \left\lceil \kappa\,\frac{Tk}{E} \right\rceil
\end{equation}
assignments. The formula is the mean load with headroom: $T$ tokens each make
$k$ requests, giving $Tk$ assignments shared by $E$ experts, so the mean is
$Tk/E$ and $\kappa$ is the slack. Worked example: $T = 4096$ tokens, $E = 8$,
$k = 2$ gives $8192$ assignments and a mean of $1024$ per expert; $\kappa =
1.5$ sets $C = 1536$. An expert that receives $2000$ requests serves its
$1536$ highest-scoring ones and drops the other $464$: those tokens leave
the layer without that expert's contribution (its weight is zeroed), and
\texttt{dropped\_fraction} records the share of removed assignments ($464$
of $8192 \approx 5.7\%$ in the example). What the cap buys is the straggler
remedy of Section~\ref{app:moe-failures}: a hard upper bound on each
expert's buffer (memory) and, in distributed training, the guarantee that no
hot expert sets the step time. The cost is silent quality loss: Pangu Ultra
MoE dropped about 6--8\% of tokens at $\kappa = 1.5$ and its dropless run
kept a lower loss, a gap that grew with the expert count
\citep{pangu2025ultra}, and MegaBlocks reached its lowest loss when no token
was dropped \citep{gale2022megablocks}. Dropless remains the default and the
recommended setting for fine-tuning, where dropped tokens can fall inside
instruction responses.

\subsection{Monitoring}
\label{app:moe-monitoring}

Each \texttt{MoELayer} exposes \texttt{last\_auxiliary\_losses}
(coefficient-weighted \texttt{moe\_load\_balancing\_loss} and
\texttt{moe\_router\_z\_loss}) and \texttt{last\_moe\_stats}. The model
aggregates them over layers through
\texttt{aggregate\_moe\_aux\_losses()} and adds the losses to the task loss;
the trainer logs the scalars to the metrics callback and to Weights \&
Biases. Read as five questions, one signal each:
\begin{itemize}[leftmargin=1.4em]
  \item \textbf{Is the load even?} MaxVio (\texttt{moe/max\_vio}) per layer:
        near $0$ healthy, $\geqslant 1$ severe (Section~\ref{app:moe-balancing});
        a value rising across steps is the collapse loop in progress.
  \item \textbf{Has the router committed?} Routing entropy and top-1 weight.
        A flat gate that merely averages its top-$k$ experts shows up as
        entropy near $\log k$ and top-1 weight near $1/k$ (no
        specialisation); a winner-takes-all gate shows entropy near $0$ and
        top-1 near $1$. Healthy training drifts from the first state towards
        committed, structured routing without reaching the second extreme.
  \item \textbf{Are the experts healthy?} The expert-norm ratio: the largest
        per-expert output norm over the layer median. Far above one means a
        rogue expert dominating the branch---the pathology the activation
        clamp contains; pinned at zero means a dead expert nobody selects
        \citep{huang2026step35}.
  \item \textbf{How hard is balancing pulling, and is the router drifting?}
        The load-balancing loss and the z-loss (coefficients already divided
        out in the logged statistics): the first should relax as balance
        improves, the second tracks logit drift.
  \item \textbf{Is capacity binding?} The dropped fraction: share of
        assignments removed by the capacity factor; it is $0$ unless
        \texttt{routing.capacity\_factor} is positive.
\end{itemize}

Five rules of thumb close the on-ramp, each a consequence of the sections
above rather than new machinery:
\begin{enumerate}[leftmargin=1.4em]
  \item Start with the defaults: dropless routing, batch-scope auxiliary
        loss at $0.01$, no bias. They are the configuration most presets
        converge to after their own deviations.
  \item If balance is the problem, reach for the bias before raising the
        auxiliary coefficient: the bias redirects traffic through the
        selection only and costs no language-model loss
        \citep{wang2024auxfree}, while a larger $\alpha$ fights the task
        loss.
  \item Watch MaxVio on the first MoE layers after any dense prefix; that is
        where collapse starts (Section~\ref{app:moe-layout}).
  \item Softmax at $k{=}1$ needs \texttt{pre\_softmax} or the router never
        learns (Section~\ref{app:moe-router}); the loader repairs it, but
        the warning is worth understanding.
  \item Budget compute as $\texttt{top\_k} \times \texttt{expert\_hidden\_size}
        + \texttt{shared}$ before enabling: shared experts and latent width
        change the per-token bill, and the cheapest fix for slow experts is
        a smaller budget, not more machinery.
\end{enumerate}

\subsection{Schema contract}
\label{app:moe-schema}

Table~\ref{tab:moe-schema} lists every \texttt{model.moe.*} key, its
default, what it controls, and the governing reference. The JSON schema is
the executable contract; runtime validation adds the cross-key rules:
$top\_k \leqslant num\_experts$, softmax at $k=1$ implies pre-softmax, the
bias requires sigmoid or sqrt-softplus scoring, expert choice is
incompatible with the balancing terms, hash layers require token choice and
forbid the bias, the latent width must lie in $(0, d)$, and the layer count
must be positive.

{\small
\begin{longtable}{@{}L{0.27\linewidth}L{0.13\linewidth}L{0.32\linewidth}L{0.20\linewidth}@{}}
\caption{MoE schema contract (\texttt{model.moe.*}).}\label{tab:moe-schema}\\
\toprule
\textbf{Key} & \textbf{Default} & \textbf{Controls} & \textbf{Reference} \\
\midrule
\endfirsthead
\toprule
\textbf{Key} & \textbf{Default} & \textbf{Controls} & \textbf{Reference} \\
\midrule
\endhead
\texttt{enabled} & \texttt{true} & routed vs.\ dense FFN & \citet{shazeer2017moe} \\
\texttt{num\_experts} & 8 & size of the routed pool & \citet{fedus2021switch} \\
\texttt{top\_k} & 2 & experts used per token & \citet{fedus2021switch} \\
\texttt{expert\_hidden\_size} & \texttt{null} & expert width (falls back to \texttt{ffn.hidden\_size}) & \citet{deepseek2024v3} \\
\texttt{router.scoring} & \texttt{softmax} & logits$\to$scores function & \citet{deepseek2024v3,deepseek2026v4} \\
\texttt{router.pre\_softmax} & \texttt{false} & softmax before the top-$k$ cut & \citet{fedus2021switch} \\
\texttt{router.renormalize} & \texttt{true} & selected weights sum to one & \citet{krajewski2025scaling} \\
\texttt{router.routed\_scaling\_factor} & 1.0 & routed-branch magnitude $g$ & \citet{deepseek2024v3,kimiteam2026k3} \\
\texttt{router.dtype} & \texttt{float32} & router compute precision & \citet{fedus2021switch,zoph2022stmoe} \\
\texttt{router.jitter\_eps} & 0.0 & input jitter strength & \citet{fedus2021switch} \\
\texttt{router.z\_loss\_coeff} & 0.0 & logit-drift penalty $\beta$ & \citet{zoph2022stmoe} \\
\texttt{routing.type} & \texttt{token\_choice} & who picks whom & \citet{zhou2022expertchoice} \\
\texttt{routing.capacity\_factor} & 0.0 & per-expert cap $\kappa$ (0: dropless) & \citet{pangu2025ultra} \\
\texttt{routing.hash\_layers} & 0 & first $N$ layers routed by table & \citet{roller2021hash,deepseek2026v4} \\
\texttt{routing.hash\_table} & \texttt{random} & table construction & \citet{roller2021hash} \\
\texttt{routing.hash\_seed} & 42 & table seed & \citet{roller2021hash} \\
\texttt{balancing.aux\_loss} & \texttt{batch} & loss scope & \citet{dai2024deepseekmoe,deepseek2024v3} \\
\texttt{balancing.aux\_loss\_coeff} & 0.01 & loss weight $\alpha$ & \citet{deepseek2024v3} \\
\texttt{balancing.bias\_update} & \texttt{none} & thermostat off / \texttt{sign} / centred & \citet{wang2024auxfree} \\
\texttt{balancing.bias\_update\_rate} & 0.001 & thermostat step $\gamma$ & \citet{wang2024auxfree} \\
\texttt{shared.count} & 0 & always-on experts & \citet{dai2024deepseekmoe} \\
\texttt{shared.hidden\_size} & \texttt{null} & shared-expert width & \citet{dai2024deepseekmoe} \\
\texttt{shared.gate} & \texttt{none} & fixed weight or learned gate & \citet{qwen2025qwen3} \\
\texttt{latent.enabled} & \texttt{false} & compress before dispatch & \citet{elango2026latentmoe} \\
\texttt{latent.dim} & \texttt{null} & latent width $\ell$ & \citet{elango2026latentmoe} \\
\texttt{latent.norm\_before\_up\_proj} & \texttt{false} & RMSNorm before $W_\uparrow$ & \citet{kimiteam2026k3} \\
\texttt{layout.dense\_layers} & 0 & dense prefix $N$ & \citet{deepseek2024v3} \\
\texttt{layout.every\_n\_layers} & 1 & MoE every $n$-th layer & \citet{meta2025llama4} \\
\texttt{experts.activation\_clamp} & 0.0 & gate-input clamp $c$ (0: off) & \citet{huang2026step35} \\
\bottomrule
\end{longtable}
}

\subsection{Named presets}
\label{app:moe-presets}

The presets in \texttt{configs/} mirror the published configurations, with
the documented single-device deviations. Table~\ref{tab:moe-presets}
summarises the MoE block of each MoE preset; the values were checked against
the models' Hugging Face \texttt{config.json} where available. Reading the
table against Section~\ref{app:moe-failures}: every recent frontier preset
answers collapse with the bias, the auxiliary loss, or both.

{\small
\begin{longtable}{@{}L{0.18\linewidth}L{0.12\linewidth}L{0.10\linewidth}L{0.14\linewidth}L{0.12\linewidth}L{0.22\linewidth}@{}}
\caption{MoE presets and their real-architecture MoE blocks.}\label{tab:moe-presets}\\
\toprule
\textbf{Preset} & \textbf{Experts / top-k} & \textbf{Expert width} & \textbf{Router} & \textbf{Shared} & \textbf{Balancing / layout} \\
\midrule
\endfirsthead
\toprule
\textbf{Preset} & \textbf{Experts / top-k} & \textbf{Expert width} & \textbf{Router} & \textbf{Shared} & \textbf{Balancing / layout} \\
\midrule
\endhead
DeepSeek-V2 & 160 / 6 & 1536 & softmax, $g{=}16$ & 2 & seq-aux $10^{-3}$; 1 dense \\
DeepSeek-V3 & 256 / 8 & 2048 & sigmoid, $g{=}2.5$ & 1 & bias + seq-aux $10^{-4}$; 3 dense \\
GLM-4.5/4.7 & 160 / 8 & 1536 & sigmoid, $g{=}2.5$ & 1 & bias; 3 dense \\
GLM-5 & 256 / 8 & 2048 & sigmoid, $g{=}2.5$ & 1 & bias; 3 dense \\
Kimi K2 & 384 / 8 & 2048 & sigmoid, $g{=}2.827$ & 1 & bias + seq-aux; 1 dense \\
Kimi K3 & 896 / 16 & 3072 & sigmoid, $g{=}1$ & 2 & centred bias; latent 3584; 1 dense \\
gpt-oss-120b & 128 / 4 & 2880 & softmax & 0 & aux, clamp 7 \\
Llama 4 Scout & 16 / 1 & 8192 & softmax, pre & 0 & every layer \\
Llama 4 Maverick & 128 / 1 & 8192 & softmax, pre & 0 & every 2nd layer \\
MiniMax-Text-01 & 32 / 2 & 9216 & sigmoid & 0 & aux $10^{-3}$ \\
MiniMax-M2 & 256 / 8 & 1536 & sigmoid & 0 & bias + aux $10^{-3}$ \\
Qwen3-Next-80B & 512 / 10 & 512 & softmax & 1, gated & aux $10^{-3}$ \\
Qwen3.8 2.4T-A95B & 512 / 10 & 2048 & softmax & 1 & aux $10^{-3}$ \\
\bottomrule
\end{longtable}
}

\paragraph{Single-device deviations.}
The runtime is CPU-first and single-device, so it implements the model-level
semantics and leaves distributed mechanisms to future work. Documented
deviations: DeepSeek-V2's device-level loss is approximated by the batch
auxiliary loss; Kimi K3's quantile/PID balancing is approximated by the
centred sign rule; Kimi K3's SiTU-GLU smooth cap is approximated by
\texttt{activation\_clamp: 4}; Llama 4's input-side score scaling is not
applied (the router weights are applied to the expert output); group-limited
routing, expert parallelism, token dispatchers, expert replicas/EPLB, and
FP8/FP4 expert execution are out of scope. DeepSeek-V4's hash layers are
supported through \texttt{routing.hash\_layers}.

\subsection{Training interaction}
\label{app:moe-training}

The expert bias is updated after the optimizer step from counts accumulated
across the step's forwards, so gradient accumulation and micro-batching do
not change its semantics. The router receives no bias gradient, and the bias
itself is a non-persistent buffer; exported checkpoints therefore contain no
training state---they remain inference-clean. Auxiliary losses are
coefficient-weighted tensors in \texttt{last\_auxiliary\_losses} and are
summed into the task loss by the trainer; they are excluded from
evaluation. The router and the experts are categorised as \texttt{router}
and \texttt{ffn} parameter blocks respectively, so per-block gradient norms
and parameter counts match the estimator taxonomy. Finally, the MoE layer is
token-wise and therefore cache-safe: the KV-cache parity suites keep the FFN
dense because a discrete top-$k$ can flip under the tiny numerical
differences between full and cached forwards---a property of discrete
routing, not of the cache.

\subsection{Related systems and future work}
\label{app:moe-future}

The appendix's design follows the MoE handbook's anatomy (router, routing,
balancing, experts, layout, capacity) and the systems literature it surveys.
Future work in this codebase: expert parallelism and all-to-all dispatchers
\citep{lepikhin2020gshard,fedus2021switch}; global-batch auxiliary losses
\citep{qiu2025demons}; quantile/PID balancing \citep{jin2026idbalancing};
dense-to-MoE upcycling and expert growth \citep{komatsuzaki2023upcycling};
frequency-balanced hash tables \citep{roller2021hash}; exact SiTU-GLU; and
shortcut-connected or parallel MoE blocks that overlap dispatch with compute
\citep{cai2024scmoe,longcat2025flash}.
